% Vietnamese translation of the complete Chapter 1; source DOI: 10.1007/978-3-030-43946-0_1.
% Source: printed 1 / PDF 17
\chapter{Mô hình hóa đa mô thức cho các hệ thống thực-ảo: Nền tảng}
\chapterauthors{Paulo Carreira, Vasco Amaral và Hans Vangheluwe}

\begin{flushleft}
\small
\textbf{Paulo Carreira}\\
Instituto Superior Técnico, Universidade de Lisbon, Bồ Đào Nha\\
Thư điện tử: \url{paulo.carreira@tecnico.ulisboa.pt}\\[0.5em]
\textbf{Vasco Amaral}\\
FCT, Universidade NOVA de Lisboa, Bồ Đào Nha\\
Thư điện tử: \url{vma@fct.unl.pt}\\[0.5em]
\textbf{Hans Vangheluwe}\\
McGill University, Canada\\
Thư điện tử: \url{hans.vangheluwe@uantwerp.be}\\[0.5em]
\copyright{} Các tác giả 2020. P. Carreira và cộng sự (chủ biên), \emph{Foundations of Multi-Paradigm Modelling for Cyber-Physical Systems}.\\
\url{https://doi.org/10.1007/978-3-030-43946-0_1}
\end{flushleft}

\section*{Tóm tắt}
Mô hình hóa và phân tích các hệ thống thực-ảo (Cyber-Physical Systems, CPS) vốn là một công việc đa ngành. Bất kỳ ai bắt đầu tham gia lĩnh vực này đều sẽ phải tìm đến một tài liệu tham khảo cung cấp nền tảng vững chắc. Mặc dù trong từng chuyên ngành, nhiều kỹ thuật đã được sử dụng như một thông lệ chuẩn, những người thực hành ở lĩnh vực khác thường vẫn xa lạ với cơ sở và cách ứng dụng các kỹ thuật đó. Nhìn chung, người thực hành có xu hướng sử dụng kỹ thuật mà họ quen thuộc nhất, bỏ qua những kỹ thuật khác cũng phù hợp với bài toán đang xét. Bản chất liên ngành của CPS đòi hỏi phải sử dụng các kỹ thuật mô hình hóa khác nhau, từ đó đặt ra nhu cầu về một hệ hình thức nền tảng chung giúp tất cả các chuyên ngành giao tiếp với nhau. Tuy nhiên, cho đến nay, chưa tồn tại một siêu hệ hình thức duy nhất như vậy để hỗ trợ nhiều chiều cạnh trong thiết kế CPS. Thực vậy, để thiết kế một CPS hiệu quả, các kỹ sư, trong vai trò người xây dựng mô hình, hoặc phải thành thạo nhiều hệ hình thức, hoặc cần có một cách tiếp cận mô hình hóa mới về căn bản. Trong chương này, chúng tôi trình bày động cơ của mô hình hóa đa mô thức cho CPS (Multi-Paradigm Modelling of CPS, MPM4CPS), giới thiệu các định nghĩa và thuật ngữ cơ bản liên quan đến mô hình hóa CPS và đa mô thức, rồi đặt nền móng cho phần còn lại của cuốn sách.

\section{Giới thiệu}
Hệ thống thực-ảo (CPS) là các hệ thống gồm những thành phần \emph{ảo} (cyber, dưới dạng các hiện thực bằng máy tính) và \emph{vật lý} (physical) [130]. Ý tưởng chung là thành phần ảo và thành phần vật lý tác động qua lại: thành phần ảo có thể khiến thành phần vật lý thay đổi trạng thái, rồi thay đổi này được phản hồi trở lại, làm thay đổi trạng thái của thành phần ảo.

Phát triển từ các khái niệm có trước, trong đó có những khái niệm thuộc cơ điện tử, hệ thống nhúng và điều khiển học, thuật ngữ ``Cyber-Physical System'' (CPS) được các tài liệu ghi nhận là do Hellen Guille đặt ra vào năm 2006 [130]. CPS thường được xem là các mạng lưới gồm những quá trình đa vật lý (cơ học, điện, sinh hóa, v.v.) và tính toán (điều khiển, xử lý tín hiệu, suy luận logic, lập kế hoạch, v.v.), thường tương tác với một môi trường đầy bất định và bất lợi, bao gồm con người và các CPS khác.

Các miền ứng dụng tiêu biểu của CPS bao gồm tiết kiệm năng lượng, điều khiển môi trường, điện tử hàng không, điều khiển hạ tầng thiết yếu (điện năng, tài nguyên nước và hệ thống thông tin liên lạc), thiết bị và hệ thống y tế có độ tin cậy cao, điều khiển và bảo đảm an toàn giao thông, hệ thống ô tô tiên tiến, điều khiển quá trình, robot phân tán (hiện diện từ xa, y tế từ xa), sản xuất và kỹ nghệ đô thị thông minh. Thiết kế CPS hiện đang thúc đẩy
% Source: printed 2 / PDF 18
đổi mới trong nhiều ngành công nghiệp, tạo ra những thị trường hoàn toàn mới. Các CPS hiệu quả hơn và rẻ hơn sẽ đem lại tác động kinh tế tích cực cho bất kỳ miền ứng dụng nào trong số này.

CPS nổi tiếng là khó thiết kế và hiện thực, chủ yếu vì chúng vượt qua ranh giới giữa các chuyên ngành và dẫn đến tương tác giữa các miền, trong những ứng dụng thường có yêu cầu an toàn trọng yếu. Thực vậy, do bản chất ứng dụng, sự cố hoặc hoạt động kém hiệu quả của CPS có thể gây ra chi phí kinh tế trực tiếp, đo lường được, làm tổn hại môi trường, hoặc thậm chí tác động trực tiếp đến con người. Vì vậy, qua nhiều năm, hoạt động kỹ nghệ đã tìm cách giải quyết những mối lo này bằng việc cải tiến các ngôn ngữ, khung và công cụ dùng trong thiết kế và phân tích CPS. Nỗ lực đó dẫn đến sự ra đời và việc áp dụng rộng rãi thiết kế dựa trên mô hình, trong đó hệ thống được thiết kế ở mức trừu tượng cao hơn, rồi một hiện thực được tạo ra bằng cách sinh tự động.

Một đặc điểm nổi bật của thiết kế CPS là bản chất đa ngành. Một nguồn gốc của tính đa ngành này nằm trong chính miền ứng dụng, chẳng hạn ngành y tế, sinh học hoặc hàng không. Một nguồn gốc khác là tính không đồng nhất của CPS, vốn gồm các bộ phận máy tính, điện tử và cơ khí. Để thiết kế một CPS, các kỹ sư thuộc nhiều chuyên ngành cần cùng nhau khảo sát các phương án thiết kế hệ thống, đi đến thống nhất, phân bổ trách nhiệm cho các phần tử phần mềm và vật lý, đồng thời phân tích những đánh đổi giữa chúng.

Bắt nguồn từ cộng đồng mô hình hóa và mô phỏng, thuật ngữ \emph{mô hình hóa đa mô thức} (Multi-Paradigm Modelling, MPM) xuất hiện từ năm 1996, khi Nhóm công tác nghiên cứu cơ bản ESPRIT 8467 của EU xây dựng một loạt hướng dẫn chính sách về mô phỏng [286], xác định nhu cầu về ``một phương pháp luận đa mô thức để biểu đạt tri thức mô hình bằng cách kết hợp những biểu diễn trừu tượng khác nhau, thay vì phát minh ra một siêu mô thức mới'', rồi sau đó đề xuất một phương pháp luận tập trung vào việc kết hợp nhiều hệ hình thức [294]. Từ năm 2004, nhiều ý tưởng về MPM đã được phát triển trong chuỗi hội thảo thường niên Computer Automated Multi-Paradigm Modelling (CAMPaM) tại Viện nghiên cứu Bellairs của Đại học McGill. Kể từ đó, MPM trở thành một lĩnh vực nghiên cứu được công nhận rộng rãi, với nhiều công trình được thực hiện và công bố, đặc biệt tại các hội thảo MPM tổ chức cùng với MoDELS.

Chương trình COST Action IC1404 MPM4CPS gần đây (\url{http://mpm4cps.eu}) hướng tới khảo sát cách sử dụng MPM để giảm độ phức tạp kỹ nghệ trong việc xây dựng ý tưởng CPS. Trong số những nỗ lực của chương trình, cộng đồng khoa học tập hợp quanh hoạt động này đã khảo sát các ngôn ngữ, kỹ thuật và công cụ hiện có thường được dùng để mô hình hóa CPS, rồi tổ chức chúng thành một bản thể học [167].

Công việc này xác định nhu cầu xây dựng nền tảng lý thuyết cho MPM và nhận diện các đặc tính ngôn ngữ hữu ích. Những đặc tính đáng chú ý nhất có thể tóm lược như sau: gần gũi với các khái niệm cốt yếu mà kỹ sư dùng để suy luận về hành vi của hệ thống vật lý, theo cách phi nhân quả về mặt tính toán (tức là không cần chỉ định sớm đâu là đầu vào và đầu ra); khả năng mô tả chính xác phép tính ở các mức chi tiết khác nhau của chiều thời gian; khả năng biểu đạt gọn gàng tính đồng thời, đồng bộ hóa và tính không tất định, cũng như suy luận về các tính chất trên mọi hành vi có thể có của một hệ thống; khả năng biểu đạt hành vi theo chế độ, có yếu tố thời gian, phản ứng và tự chủ, đồng thời tổng hợp mã; sự phù hợp để mô hình hóa cạnh tranh tài nguyên dùng chung, dẫn đến xếp hàng, làm cơ sở cho phân tích hiệu năng định lượng; sự phù hợp để hợp thành kiến trúc một cách dễ dàng và đúng đắn; khả năng biểu đạt luồng công việc ở mức trừu tượng cao; và cuối cùng là đặc tính cấp cao về tính mô-đun, hỗ trợ tái sử dụng, kiểm chứng tính chất theo phép hợp thành và tích hợp các thành phần hộp đen như các đơn vị đồng mô phỏng. Phạm vi của các kỹ thuật này rất rộng.

Kết quả của các nỗ lực nêu trên trong cộng đồng MPM đã kết tinh thành một cuốn sách tổng hợp nhất quán những tri thức đã được xác lập về nền tảng và các hệ hình thức để mô hình hóa CPS, đặc biệt chú trọng các công cụ và kỹ thuật đa mô thức.

\section{Tìm hiểu các hệ thống thực-ảo}
Giờ đây, chúng tôi chuyển sang thảo luận những lớp hệ thống nào có thể được xem là CPS, các tính chất của chúng là gì và những nguồn gây ra độ phức tạp kỹ nghệ nào đáng được giải quyết bằng cách tiếp cận gọi là MPM. Trước khi bắt đầu thảo luận, cần làm rõ một số khái niệm sơ bộ.

% Source: printed 3 / PDF 19
\subsection{Hệ thống và các mô hình của chúng}
\emph{Hệ thống} là một khái niệm cơ bản được dùng trong nhiều chuyên ngành. Thuật ngữ này khái niệm hóa một thực thể vật lý như hệ sinh thái, sinh vật, máy móc, hoặc một thực thể thuần túy \emph{trừu tượng}. Trường hợp sau đề cập đến các quá trình, quy tắc (như một hệ thống kinh tế--xã hội), hoặc các mô hình toán học. Trong mọi trường hợp, một hệ thống được hiểu là (i) một tập hợp \emph{thành phần}, (ii) có thể nhận diện như một \emph{tổng thể}, (iii) \emph{hợp tác} (cũng gọi là tương tác) với nhau để (iv) thực hiện một \emph{chức năng} cụ thể. Cần phân biệt các hệ thống thực tế (vật lý hoặc trừu tượng) với hiểu biết của con người về chúng. Việc hình thức hóa hiểu biết này, vốn thường có giới hạn, được gọi là một \emph{mô hình} [308].

Theo Lý thuyết hệ thống, có thể hiểu một hệ thống theo hai cách. Một cách là tiếp cận \emph{hộp đen}, nhằm mô hình hóa \emph{hành vi bên ngoài} của hệ thống theo cách nó tương tác với môi trường. Ở đây, hành vi của hệ thống được xem là quan hệ giữa sự tiến triển của lịch sử các \emph{biểu hiện} (cũng gọi là \emph{đầu ra}) và lịch sử các \emph{tác động kích thích} (cũng gọi là \emph{đầu vào}). Cách hiểu thứ hai là tiếp cận \emph{hộp trắng}, nhằm hiểu \emph{cấu trúc bên trong} của hệ thống theo các \emph{thành phần và kết nối} mà qua đó tương tác diễn ra.

Cả tiếp cận hộp đen lẫn hộp trắng đều hình thức hóa tri thức về hành vi hệ thống bằng các mô hình\footnote{Ý nghĩa chính xác của thuật ngữ ``mô hình'' được bàn luận ở Mục 1.3.1.}, là những biểu diễn đơn giản hóa của hệ thống. Giá trị của các biểu diễn sử dụng được nằm ở chỗ chúng là phương tiện hiệu quả về mặt nhận thức để suy luận về hệ thống thực tế (hoặc các hệ thống) mà chúng đại diện. Suy luận ở đây có nghĩa là phân tích các tính chất của hệ thống hoặc dự đoán hành vi tương lai của hệ thống. Vì thế, một hệ thống có thể có nhiều mô hình phục vụ những yêu cầu suy luận khác nhau. Về mặt lịch sử, Lý thuyết điều khiển sử dụng tiếp cận hộp đen, trong đó đầu ra được mô hình hóa theo đầu vào bằng các phương trình vi phân; Khoa học máy tính truyền thống sử dụng tiếp cận hộp trắng, mô hình hóa cấu trúc bên trong của hệ thống bằng các biểu đồ đối tượng.

Ý tưởng chủ đạo của việc tìm hiểu cấu trúc bên trong là: các thành phần được liên kết theo một cách bố trí kết nối nhất định sẽ thể hiện hành vi bên ngoài cụ thể. Hơn nữa, hành vi tổng thể có thể được suy ra từ (i) hành vi đã biết của các thành phần và (ii) đặc điểm của các kết nối. Kết nối là những dòng thông tin hoặc năng lượng trừu tượng, liên kết các thành phần thông qua những điểm tương tác được chỉ định trước, gọi là \emph{cổng}. Một số cổng thành phần có hướng, nên có thể nói đến \emph{cổng đầu vào} (đầu vào) và \emph{cổng đầu ra} (đầu ra). Những cổng khác không có hướng, theo nghĩa chúng mô hình hóa sự trao đổi chứ không nhất thiết là một dòng chảy (chẳng hạn một bộ ghép nhiệt).

Một thành phần cũng có thể được xem là một hệ thống với các thành phần và kết nối của riêng nó. Khả năng tiếp tục hợp thành hệ thống từ những hệ thống đã được xây dựng trước đó được gọi là \emph{xây dựng phân cấp}. Các thành phần cũng thường được mô hình hóa với một \emph{trạng thái bên trong}, phản ánh (một phần và đôi khi không chính xác) những gì hệ thống biết về môi trường xung quanh, cùng một cơ chế \emph{chuyển trạng thái}, thường gọi là \emph{hàm chuyển}, có khả năng tạo ra trạng thái mới từ trạng thái trước đó khi nhận một số đầu vào và tạo ra một số đầu ra.

\subsection{Các loại hệ thống}
Đầu vào, đầu ra và trạng thái của các thành phần được mô hình hóa bằng các \emph{biến} nhận giá trị từ các tập giá trị tương ứng. Các biến này được gọi là \emph{liên tục} hoặc \emph{rời rạc} nếu tập giá trị của chúng, theo nghĩa toán học, lần lượt là tập trù mật hoặc tập rời rạc. Hành vi của hệ thống, được mô hình hóa theo trạng thái bên trong và đầu ra, tiến triển từ một trạng thái trước đó theo một khái niệm thời gian nào đó. Thời gian cũng có thể được hiểu là liên tục hoặc rời rạc. Trong \emph{thời gian liên tục}, có thể suy ra trạng thái và đầu ra mới của hệ thống với một khoảng thời gian nhỏ tùy ý; còn trong \emph{thời gian rời rạc}, chỉ có thể suy ra trạng thái và đầu ra mới tại những khoảng thời gian định trước, hoặc khi một sự kiện nhất định xảy ra.

Có thể phân loại hệ thống thành các loại khác nhau tùy theo bản chất của đầu vào, đầu ra, biến trạng thái, hàm chuyển trạng thái và khái niệm thời gian. Một cách phân loại được chấp nhận rộng rãi là như sau [308]:
\begin{itemize}
\item \textbf{Hệ thống tĩnh và hệ thống động.} Một hệ thống được gọi là \emph{tĩnh} nếu đầu ra chỉ phụ thuộc vào đầu vào hiện tại. Nếu đầu ra của hệ thống phụ thuộc vào lịch sử các đầu vào trước đó, hệ thống được gọi là \emph{động}.
% Source: printed 4 / PDF 20
\item \textbf{Nhân quả và phi nhân quả (không nhân quả).} Khi giá trị đầu ra của hệ thống độc lập với các giá trị đầu vào trong tương lai, hệ thống được gọi là \emph{hệ thống nhân quả}; khi các giá trị đầu ra của hệ thống phụ thuộc vào giá trị đầu vào tại bất kỳ thời điểm nào, hệ thống được gọi là \emph{hệ thống phi nhân quả}.
\item \textbf{Tuyến tính và phi tuyến.} Một hệ thống được gọi là \emph{tuyến tính} (tương ứng, \emph{phi tuyến}) nếu sự thay đổi ở đầu ra tỷ lệ (tương ứng, không tỷ lệ) với sự thay đổi của đầu vào.
\item \textbf{Trạng thái rời rạc và trạng thái liên tục.} Những hệ thống trong đó các biến trạng thái chỉ thay đổi tại một tập rời rạc các thời điểm được gọi là hệ thống \emph{trạng thái rời rạc}; những hệ thống trong đó các biến trạng thái thay đổi liên tục theo thời gian (ví dụ một bể nước đang được đổ đầy) được gọi là hệ thống \emph{trạng thái liên tục}.
\item \textbf{Thời gian rời rạc và thời gian liên tục.} Một hệ thống được gọi là \emph{rời rạc}, đối lập với \emph{liên tục}, nếu nó có số trạng thái đếm được.
\begin{quote}\small
\textbf{Ghi chú bản dịch (C1-1).} Phân loại thời gian rời rạc hay liên tục phụ thuộc vào miền thời gian, không phụ thuộc vào số trạng thái. Chẳng hạn, mô hình $x_{n+1}=0.5x_n$ với $n\in\mathbb{N}$ và $x_n\in\mathbb{R}$ có thời gian rời rạc nhưng không gian trạng thái không đếm được. Câu trên được giữ theo bản gốc.
\end{quote}
\item \textbf{Dẫn động theo thời gian và dẫn động theo sự kiện.} Một hệ thống \emph{dẫn động theo thời gian} thay đổi trạng thái theo một thời gian vật lý đều đặn. Trong khi đó, một hệ thống \emph{dẫn động theo sự kiện} thay đổi trạng thái để phản ứng với sự xuất hiện của các sự kiện rời rạc không đồng bộ (và không thay đổi trạng thái giữa các lần xuất hiện sự kiện).
\item \textbf{Biến thiên theo thời gian và bất biến theo thời gian.} Một hệ thống mà trong đó một số biến trạng thái thay đổi theo thời gian, khiến hệ thống phản ứng khác nhau với cùng một đầu vào tại các thời điểm khác nhau, được gọi là \emph{biến thiên theo thời gian}; một hệ thống cho cùng một đầu ra với một đầu vào nhất định tại những thời điểm khác nhau được gọi là \emph{bất biến theo thời gian}.
\end{itemize}

Một hệ thống động thể hiện cả hành vi thời gian liên tục lẫn thời gian rời rạc được gọi là \emph{hệ thống lai}. Nghiên cứu hệ thống lai rất quan trọng vì chúng thường xuất hiện khi hợp thành các thành phần rời rạc với thành phần liên tục, một đặc điểm điển hình của hệ thống thực-ảo.

\subsection{Hệ thống thực-ảo là gì?}
Thành phần ảo được giả định là \emph{điều khiển} thành phần vật lý, theo nghĩa thành phần ảo có một mức ``trí thông minh'' nào đó, hoặc ít nhất một chiến lược để dẫn dắt phần vật lý đạt được một mục tiêu quan sát được đã định trước. Chiều ngược lại không đúng, tức là phần vật lý không nhằm dẫn dắt phần ảo đạt một mục tiêu định trước nào đó. Cách phát biểu này đặt CPS trong phạm vi của Lý thuyết hệ thống điều khiển. Mặc dù Lý thuyết hệ thống điều khiển có tính bao quát, CPS không phải là một trường hợp riêng của hệ thống điều khiển.

Trong một \emph{hệ thống điều khiển}, một thành phần gọi là \emph{bộ điều khiển} hiện thực một mô hình điều khiển tác động lên thành phần môi trường vật lý, gọi là \emph{đối tượng điều khiển} (plant), thông qua một \emph{tác động điều khiển}. Bộ điều khiển biết giá trị mong muốn của một biến và giá trị hiện tại của biến đó trong thành phần vật lý. Sau đó, bộ điều khiển tạo ra một chuỗi tác động điều khiển để hiệu chỉnh một biến (đo được từ hệ thống vật lý) về giá trị mong muốn (mục tiêu), tức là giảm thiểu khoảng cách theo thời gian. Cách bố trí này được gọi là \emph{hệ thống điều khiển}. Vì vậy, có thể đặt câu hỏi: ``Điều gì phân biệt một hệ thống điều khiển với một CPS?'' Ngoài chính ý tưởng về một thành phần điều khiển ảo---một thuật toán máy tính rời rạc điều khiển hiện tượng vật lý liên tục---dường như không có một đặc điểm riêng lẻ nào tự nó định nghĩa được CPS. Tuy nhiên, người ta nhìn chung chấp nhận rằng hệ thống thực-ảo gồm rất nhiều thành phần tương tác với nhau và thể hiện một số đặc điểm thường gặp, phân biệt chúng với các hệ thống điều khiển kinh điển. Cụ thể là:
\begin{itemize}
\item \textbf{Các thành phần ``ảo'' rộng khắp}, mã hóa logic điều khiển và điều khiển giám sát phức tạp. Thông thường, chúng có nhiều thành phần ảo con hỗ trợ phối hợp hành động phức tạp và đòi hỏi xử lý khối lượng dữ liệu lịch sử rất lớn.
\item \textbf{Quy mô hoạt động rất lớn}, vượt quá vài triệu phần tử (đôi khi không đồng nhất), kéo theo độ phức tạp vốn có của các cấu trúc phân cấp và tương tác. Ví dụ gồm lưới điện thông minh, đô thị thông minh, máy dò hạt trong vật lý hạt, cùng những hệ thống khác.
\item \textbf{Bản chất lai rời rạc--liên tục}: rất nhiều thành phần rời rạc (đặc biệt là các thành phần ``ảo'') được kết nối với các thành phần vật lý (có bản chất liên tục), tạo thành một hệ thống lai. Ngoài ra, nhiều thành phần trong số này khá không đồng nhất về loại hệ thống (Mục 1.2.2).
\item \textbf{Tích hợp với nhiều hệ thống bên ngoài}, bằng cách xử lý dữ liệu và sự kiện ở các định dạng khác nhau từ nhiều hệ thống có băng thông và mức bảo đảm chuyển giao thông điệp khác nhau.
\item \textbf{Kết nối mạng sâu rộng và có cấu trúc phân cấp}, kết nối nhiều thành phần, thường thông qua mạng số với những bus và giao thức truyền thông khác nhau.
\item \textbf{Có khả năng thích ứng}: hệ thống phải điều chỉnh hành vi theo những mẫu mà ở thời điểm thiết kế chưa thể tính đến.
% Source: printed 5 / PDF 21
\item \textbf{Có con người trong vòng lặp}, với những cơ chế dành riêng để con người sử dụng đầu ra và cung cấp đầu vào. Con người thường được mô hình hóa như các thành phần có những yêu cầu và giả định cụ thể về hành vi.
\end{itemize}

Sự phổ biến của IoT với số lượng lớn thiết bị kết nối đã tạo ra một làn sóng quan tâm đến CPS. Tuy nhiên, thuật ngữ này phần nào bị lạm dụng và nay trở thành tên gọi bao trùm mọi hệ thống tương tác với môi trường vật lý. Chẳng hạn, trong thực tiễn hiện nay, một người phát triển bộ ghi nhiệt độ cũng có thể tuyên bố thiết bị của mình là CPS. Điều này đặt ra câu hỏi: \emph{Đâu là những yêu cầu tối thiểu để một hệ thống được xem là CPS?} Một nhận xét then chốt là việc có những thành phần ảo và vật lý nhận diện được chưa đủ để một hệ thống được xem là thực-ảo; các thành phần ảo và vật lý phải tác động đến nhau\footnote{Về mặt mô hình hóa, điều này có nghĩa là chúng phải được liên kết qua ít nhất một kết nối.}. Một câu hỏi hợp lý khác là liệu quan hệ này có phải là tác động qua lại hay không, hay thành phần ảo có thể không chịu tác động của thành phần vật lý. Rõ ràng, vì phần ảo điều khiển phần vật lý, phần ảo phải có phương tiện để tác động lên phần vật lý. Do đó, một hệ thống mà thành phần ảo không tác động đến thành phần vật lý thì không phải CPS.

CPS cũng được xây dựng trên nền tảng \emph{hệ thống nhúng}, là những hệ thống độc lập, tích hợp các phần tử logic điều khiển và tương tác với thế giới thực. Một hệ thống nhúng thường là một thiết bị đơn lẻ, trong khi CPS bao gồm nhiều hệ thống thành phần. Hơn nữa, các hệ thống nhúng được thiết kế chuyên biệt để thực hiện một số lượng nhiệm vụ hạn chế, thường với nguồn lực hạn chế. Ngược lại, một CPS hoạt động ở quy mô lớn hơn nhiều, có thể bao gồm nhiều hệ thống nhúng hoặc các phần tử CPS khác, kể cả con người và các hệ thống xã hội--kỹ thuật.

Từ những điều trên, có thể thấy rõ rằng việc có một thành phần ảo và một thành phần vật lý không phải điều kiện đủ để một hệ thống được xem là hệ thống thực-ảo. Hệ thống ghi nhiệt độ trong ví dụ, vốn chỉ đọc từ hệ thống vật lý, không phải CPS dù có cả thành phần ảo lẫn thành phần vật lý.

\subsection{Các nguồn gây ra độ phức tạp kỹ nghệ}
Hệ thống thực-ảo khó xây dựng vì độ phức tạp vốn có của các bài toán mà chúng giải quyết: phối hợp hành động để tối ưu hóa nhiều mục tiêu (có thể mâu thuẫn) trong những hệ thống có số lượng rất lớn cảm biến và cơ cấu chấp hành. Một nguồn phức tạp khác mang tính phụ phát nhiều hơn và liên quan đến tính không đồng nhất của các thành phần.

Một nguồn phức tạp ban đầu là các thành phần ảo có bản chất rời rạc nhưng phải hoạt động theo thời gian; vì vậy, chúng phải suy luận về thời gian. Các khái niệm \emph{thời lượng}, \emph{hạn chót} và \emph{tính đồng thời về thời điểm} phải được xử lý và mô hình hóa một cách tường minh. Tuy nhiên, các ngôn ngữ lập trình trừu tượng hóa bỏ khái niệm thời gian và cung cấp rất ít hoặc không cung cấp hỗ trợ nào cho thời gian\footnote{Ngữ nghĩa chặt chẽ về thời gian không có trong các ngôn ngữ máy tính tiêu chuẩn.}. Hành vi thời gian được thực hiện thông qua phần cứng định thời chuyên dụng, các thủ tục điều khiển ngắt và các thư viện bộ định thời---theo lời Eduard Lee, ``lập trình viên phải bước ra ngoài sự trừu tượng của lập trình để chỉ định hành vi thời gian''. Thay vì được mô hình hóa tường minh, hành vi thời gian bị che khuất và vùi dưới độ phức tạp của việc phối hợp những cơ chế này, khiến rất khó bảo đảm tính đúng đắn khi hợp thành các thành phần ảo và vật lý. Vấn đề càng trầm trọng khi thêm nhiều thành phần vào hệ thống.

Không chỉ các thành phần rời rạc thường thể hiện hành vi thời gian không đúng, việc ghép nối các thành phần rời rạc và liên tục còn đặt ra một thách thức kỹ nghệ bất ngờ khác. Hợp thành một mô hình tất định của phần ảo với một mô hình tất định của phần vật lý tạo ra một mô hình không tất định rất khó phân tích. Để hiểu nguyên nhân, hãy xét rằng trên thực tế không thể bảo đảm các thành phần hiện thực hai mô hình sẽ hoàn toàn khớp nhau về thời gian. Điều này đặc biệt đúng với cơ chế mà chúng dùng để giao tiếp và đồng bộ hóa, vì cơ chế đó hoạt động theo những giả định (ràng buộc) nhất định về thời gian.

Việc hiện thực CPS trong thực tế chủ yếu dựa trên thành phần. Các thành phần này đến từ nhiều nhà cung cấp và sử dụng nhiều công nghệ; vì vậy, chúng thường có giao thức truyền thông, thời gian đáp ứng, dung sai và điều kiện hoạt động khác nhau. Một mặt của vấn đề là một số ràng buộc trong số này không được biết hoặc tính đến từ đầu. Mặt khác là trong quá trình xây dựng hệ thống, các ràng buộc này không được mô hình hóa và xử lý tường minh. Nhiều độ phức tạp phụ phát xuất hiện khi cố gắng lắp ráp các thành phần có mức biến thiên như vậy, đặc biệt vì một CPS điển hình gồm rất nhiều thành phần.

% Source: printed 6 / PDF 22
\sourcefigure{1.1}{22}{Minh họa câu chuyện sáu người mù. Mỗi người, từ góc nhìn của mình, hiểu con voi như một vật khác nhau.}

CPS còn có những yêu cầu nghiêm ngặt không thể nới lỏng. Ngoài những bảo đảm thông thường về tính đúng đắn chức năng áp dụng cho mọi đối tượng được tạo ra bằng kỹ nghệ, CPS đặc biệt được biết đến với số lượng lớn các yêu cầu ngoài chức năng (cũng gọi là ``phi chức năng''), bao quát độ tin cậy, hiệu năng, an toàn, an ninh và những vấn đề khác. Điều cần lưu ý là các yêu cầu ngoài chức năng đặt ra những ràng buộc phức tạp đối với thiết kế thành phần định thời và hành vi thời gian. Hơn nữa, vì các hệ thống này có tính trọng yếu, chúng thường phải trải qua những quy trình đánh giá đủ điều kiện/kiểm thử nghiêm ngặt. Khi cần thay đổi hệ thống, người ta thường ngại thực hiện vì chi phí liên quan, khiến CPS khó thích ứng với các yêu cầu thay đổi. Cần có một phương tiện đáng tin cậy để bảo đảm một số khía cạnh của hệ thống không bị tác động, và vì vậy không cần được kiểm thử lại.

Các hiện thực thực tế của thành phần đôi khi tương tác theo những cách chưa được thiết kế từ đầu. Các thành phần vật lý thường gây nhiễu lẫn nhau do nhiễu điện từ hoặc nhiệt. Hành vi nảy sinh là hành vi phức tạp xuất hiện từ tương tác giữa các thành phần mà trước đó không được dự kiến. Khi phát triển CPS, điều quan trọng là phải có phương tiện khảo sát các phương án thiết kế khác nhau.

\section{Mô hình hóa một hệ thống thực-ảo}
Quy trình kỹ nghệ CPS đòi hỏi sự tham gia của nhiều chuyên ngành. Mỗi chuyên ngành thường tạo ra một thiết kế (cũng gọi là một \emph{khung nhìn} hoặc \emph{cách hiểu}) về hệ thống cho mục đích riêng, dưới dạng một mô hình. Mô hình được tạo ra bằng các phép trừu tượng hóa, các phương pháp phân rã theo kinh nghiệm và những công cụ phân tích đặc trưng của từng chuyên ngành. Mỗi chuyên ngành cũng đem đến một khối tri thức giúp những người chịu trách nhiệm mô hình hóa thực tại đánh giá có phê phán tính đúng đắn và vững chắc của mô hình đang tạo ra. Theo nghĩa này, cách tiếp cận từng phần với các mô hình riêng biệt được tổ chức theo chuyên ngành có hiệu quả. Một lý do khác để có các mô hình khác nhau là ngay trong cùng một chuyên ngành, chúng vẫn cần trả lời những câu hỏi khác nhau. Mô hình còn phải được tạo ra ở các mức chi tiết khác nhau; do đó, người xây dựng mô hình phải tìm được sự cân bằng và nắm bắt đúng những yếu tố cần thiết để tạo một mô hình phù hợp với câu hỏi đang nghiên cứu.

Mô hình hóa CPS vốn được thể hiện qua nhiều khung nhìn của hệ thống (phần lớn tuân theo nguyên tắc tách biệt các mối quan tâm). Như trong câu chuyện cổ về sáu người mù (xem Hình 1.1), không một khung nhìn đơn lẻ nào (cũng như hệ hình thức mô hình hóa tương ứng) có thể mô hình hóa mọi khía cạnh của một hệ thống. Tương tự, trong kỹ nghệ CPS, kết quả là các mô hình phản ánh những khung nhìn khác nhau về bài toán, được biểu đạt bằng nhiều ký pháp. Một vấn đề của kỹ sư CPS, cũng như của những người mù, là ``làm thế nào tích hợp tri thức để hình thành một mô hình gần với thực tại hơn?'' Một cách tự nhiên, mô hình hóa CPS đòi hỏi cách tiếp cận xuyên ngành, hợp nhất các mô hình khác nhau thành một sự trừu tượng hóa thống nhất về thực tại. Do đó, một thách thức thiết yếu của CPS là làm thế nào
% Source: printed 7 / PDF 23
kết hợp các phép trừu tượng hóa của những chuyên ngành kỹ nghệ khác nhau với các mô hình cho quá trình vật lý, bao gồm phương trình vi phân, quá trình ngẫu nhiên và những mô hình khác.

\sourcefigure{1.2}{23}{Các mô hình của một máy bay: bên trái là mô hình CAD để nghiên cứu thể tích; bên phải là mô hình để nghiên cứu trọng lực.}

Hiện chưa có một cách tiếp cận thiết kế và mô hình hóa tiêu chuẩn để tích hợp các mô hình do các chuyên ngành khác nhau trong CPS tạo ra. Thực vậy, tri thức được nắm bắt trong các mô hình chỉ được kết hợp khi lắp ráp những hiện thực vật lý (nguyên mẫu) để đánh giá một số tính chất liên quan của một CPS nhất định. Một cách tiếp cận tinh vi hơn là hoàn toàn tránh tạo ra hiện thực nguyên mẫu và thay vào đó mô phỏng hoặc đồng mô phỏng các mô hình. Cách tiếp cận này tổng quát hơn vì cho phép kiểm chứng tính chất độc lập với những đặc thù của nguyên mẫu đã được hiện thực.

Trước đó, chúng tôi đã dùng thuật ngữ \emph{mô hình} với một quan niệm ngầm định về mô hình hóa dựa trên \emph{trừu tượng hóa}: quá trình loại bỏ chi tiết khi nghiên cứu đối tượng hoặc hệ thống để tập trung vào những chi tiết quan trọng hơn. Tuy nhiên, không phải mọi phép trừu tượng hóa hoặc biểu diễn có thể nghĩ ra về một hệ thống đều được xem là mô hình. Thực vậy, cần hiểu chính xác hơn các khái niệm này. Hãy xem xét chúng kỹ hơn.

\subsection{Mô hình là gì?}
Mô hình của một hệ thống là một sự trừu tượng hóa (một biểu diễn) để đưa ra dự đoán hoặc suy luận [176] về một thực tại. Cụ thể hơn, thực tại này là một \emph{hệ thống đang nghiên cứu} (system under study, SUS), mà chúng ta muốn tìm hiểu các quy tắc chi phối và tính chất của nó trong bối cảnh một \emph{khung thực nghiệm} nhất định. Khung thực nghiệm biểu thị tập hợp giới hạn những hoàn cảnh mà trong đó một hệ thống được quan sát hoặc đưa vào thực nghiệm.

Theo Stachowiak [256], một mô hình cần có ba tính chất chính. Tính chất thứ nhất là \emph{đặc tính ánh xạ}: để được gọi là mô hình, nó phải trung thành với hoặc dựa trên một nguyên bản đang tồn tại, hoặc đơn giản là sự hình thức hóa một ý tưởng sẽ được hiện thực vào một thời điểm trong tương lai. Tính chất thứ hai là \emph{đặc tính lược giảm}. Theo đó, một biểu diễn không phải mô hình nếu nó không loại bỏ chi tiết không cần thiết và lựa chọn những tính chất của nguyên bản hữu ích cho mục đích của mô hình đang xét. Tính chất thứ ba là \emph{đặc tính thực dụng}. Theo tính chất cuối này, mô hình cần có thể được sử dụng thay cho nguyên bản đối với một mục đích nào đó. Những đặc điểm quan trọng khác của mô hình, dù không phải đặc điểm nền tảng, là \emph{có mục đích} [243], \emph{dễ hiểu} và \emph{hiệu quả về chi phí}.

Để làm rõ hơn ý tưởng về mô hình, hãy minh họa những tính chất này bằng một ví dụ đơn giản. Xét một máy bay, trong đó chúng ta tập trung nghiên cứu riêng các tính chất khối lượng và trọng lực. Cho mục đích này, có thể nói một cốc nước là một phép xấp xỉ thô của máy bay (và được chọn làm biểu diễn vật lý của SUS). Khi rơi, cốc mất một khoảng thời gian để chạm sàn, rồi vỡ. Có thể xem cốc là mô hình của máy bay không? Để đánh giá tác động của trọng lực (phạm vi của chúng ta), có (i) đặc tính ánh xạ, vì đối tượng (cốc) đại diện cho nguyên bản (máy bay); (ii) đặc tính lược giảm, vì mọi chi tiết không cần thiết như hình dạng, khí động học, kiến trúc và những chi tiết khác đều được loại bỏ; và (iii) tính thực dụng, vì có thể dùng cốc thay cho máy bay thật. Cuối cùng, nó (iv) có mục đích, vì nhằm nghiên cứu trọng lực; (v) dễ hiểu, vì là một biểu diễn đơn giản mà mọi người đều hiểu; và (iv) hiệu quả về chi phí: để nghiên cứu và thay thế vật thật trong khi rơi, nó rẻ hơn nhiều bậc độ lớn và chấp nhận được về mặt đạo đức. Một ví dụ khác về mô hình cho cùng nghiên cứu tình huống có thể là bản vẽ CAD
% Source: printed 8 / PDF 24
của máy bay nhằm nghiên cứu thể tích của nó, chẳng hạn nghiên cứu cách nó vừa với một nhà chứa máy bay, hoặc thể tích bên trong để lựa chọn hệ thống thông gió phù hợp.

\subsection{Nhiều hệ hình thức trong CPS}
Như đã thấy ở trên, quy trình kỹ nghệ CPS tạo ra một tập hợp các mô hình. Mô hình là những sự trừu tượng hóa của một hệ thống (một thực tại) có các tính chất đáng nghiên cứu. Cũng đã biết rằng mô hình, đặc biệt là mô hình thuộc những chuyên ngành khác nhau, được biểu đạt bằng các hệ hình thức mô hình hóa khác nhau tương ứng. Một \emph{hệ hình thức mô hình hóa} (modelling formalism) là một ngôn ngữ có cú pháp và ngữ nghĩa hình thức. Ý nghĩa thông thường của ``hình thức'' là được định nghĩa chính xác và không nhập nhằng bằng toán học, dưới dạng phương trình vi phân, ô-tô-mát trạng thái hữu hạn, Statecharts (biểu đồ trạng thái), mạng Petri và những hệ hình thức khác. Các hệ hình thức khác nhau tồn tại vì chúng ngắn gọn hơn và cho phép trả lời hiệu quả các lớp câu hỏi khác nhau. Thực vậy, không một hệ hình thức đơn lẻ nào có thể được dùng để mô hình hóa mọi khía cạnh của một hệ thống, vì hệ hình thức cần dùng phụ thuộc vào bản chất của bài toán cần giải quyết.

Các hệ hình thức cũng cho phép thao tác và viết lại mô hình, hoặc các phần của mô hình, thành những mô hình khác tương đương với nguyên bản (đối với một ngữ nghĩa cho trước), mà các hiện thực của chúng có những tính chất mong muốn hơn. Chẳng hạn, ít dư thừa hơn, gọn hơn, nhanh hơn hoặc tiêu thụ ít năng lượng hơn.

Để vượt qua độ phức tạp của bài toán, một thông lệ mô hình hóa phổ biến là mô tả các mô hình của cùng một thực tại ở những mức trừu tượng khác nhau (đôi khi sử dụng các hệ hình thức khác nhau tương ứng). Các mô hình được biểu đạt ở những mức trừu tượng khác nhau được liên kết với nhau thông qua các ánh xạ bảo toàn cấu trúc. Thực vậy, một vấn đề bao trùm khi dùng các hệ hình thức khác nhau là hợp nhất các mô hình của cùng một hệ thống thông qua những ánh xạ này. Cần có một khái niệm về \emph{tính nhất quán} giữa chúng.

Những hệ hình thức này được dùng để mô hình hóa một hệ thống đang tương tác với môi trường, kiến trúc và các thành phần của hệ thống, ở những mức chi tiết, xấp xỉ và trừu tượng khác nhau, từ những góc nhìn khác nhau, cũng như các nền tảng mà thành phần phần mềm của hệ thống sẽ được triển khai trên đó. Việc tích hợp các mô hình tạo ra theo những hệ hình thức khác nhau được thực hiện bằng cách ánh xạ (biên dịch) mô hình sang các \emph{siêu hệ hình thức} ở mức thấp hơn, tích hợp các miền khác nhau như đồ thị liên kết năng lượng (bond graph), hoặc các hệ hình thức khác tích hợp cấu trúc mô hình hóa rời rạc và liên tục như DEVS.

Để hỗ trợ thiết kế CPS, cần sử dụng nhiều hệ hình thức mô hình hóa được chọn theo những đặc tính riêng về suy luận và phân tích, thay vì một siêu hệ hình thức duy nhất. Những đặc tính này khiến mỗi hệ hình thức phù hợp nhất với một mục tiêu (hoặc mục tiêu con) cụ thể trong thiết kế CPS. Về mặt thực hành, kỹ sư trong vai trò người xây dựng mô hình cũng cần biết các chiến lược (được hình thức hóa thành \emph{quy trình}) để mô tả thực tại theo hệ hình thức. Hệ hình thức cùng với những quy trình và ràng buộc này tạo thành cái được gọi là một \emph{mô thức mô hình hóa}---đối tượng nghiên cứu của mô hình hóa đa mô thức [286].

\section{Mô hình hóa đa mô thức cho CPS}
Mô hình hóa đa mô thức (MPM) được công nhận là một cách tiếp cận mạnh (bản thân nó cũng là một mô thức), có thể hữu ích trong việc thiết kế, giao tiếp và suy luận về CPS, vốn nổi tiếng phức tạp vì vượt qua ranh giới giữa các chuyên ngành và có tương tác liên miền.

Để phát triển một CPS, người quản lý dự án và kỹ sư cần chọn các ngôn ngữ phát triển, vòng đời phần mềm và ``giao diện'' phù hợp nhất để đặc tả những khung nhìn, thành phần và tương tác khác nhau của hệ thống với ít ``độ phức tạp phụ phát'' (accidental complexity) [52] nhất có thể. Chẳng hạn, khi biết rằng yêu cầu hệ thống/phần mềm có thể thay đổi thường xuyên trong quá trình dự án, lựa chọn một quy trình phát triển Agile có thể giúp ứng phó với tiến hóa và thay đổi. Nếu hành vi của hệ thống đòi hỏi các thao tác được kích hoạt khi có dữ liệu, tương tự hệ thống phản ứng, các ngôn ngữ luồng dữ liệu có thể giúp đặc tả chính xác những phần trọng yếu nhất của hành vi phần mềm, để có thể phân tích thời gian.

% Source: printed 9 / PDF 25
\subsection{Mô thức là gì?}
Trong Khoa học máy tính, các ngôn ngữ lập trình đa dụng (general-purpose programming languages, GPL) có thể được phân loại theo một hoặc nhiều mô thức mà chúng hỗ trợ. Chẳng hạn, Eiffel là ngôn ngữ hướng đối tượng và hỗ trợ mô thức thiết kế dựa trên hợp đồng, Prolog là ngôn ngữ khai báo, còn Lisp là ngôn ngữ hàm. Mô thức đặc trưng cho các cấu trúc cú pháp, ngữ nghĩa nền tảng và các nguyên tắc chi phối những GPL này. Cụ thể, hướng đối tượng có bản chất mệnh lệnh và buộc ta nhìn thế giới dưới dạng các lớp và đối tượng giao tiếp với nhau, trong khi phong cách khai báo dựa trên thay thế và viết lại hạng thức. Do đó, một câu lệnh Eiffel có rất ít điểm chung với một mệnh đề Prolog vì mỗi ngôn ngữ hỗ trợ một cách nhìn rất khác. Một mô thức lập trình trực tiếp chuyển thành những khái niệm khác nhau được mã hóa trong định nghĩa cú pháp của GPL (được gọi là một \emph{siêu mô hình} trong lĩnh vực kỹ nghệ hướng mô hình, Model-Driven Engineering). Một cách tự nhiên, ý tưởng kết hợp nhiều mô thức ở cấp GPL dẫn đến những ngôn ngữ lập trình giàu sức biểu đạt và mạnh hơn như Java (mang tính mệnh lệnh, hướng đối tượng, đồng thời, thời gian thực và gần đây là tính hàm) và Maude (mang tính khai báo, hướng đối tượng, đồng thời và thời gian thực).

Vậy \emph{mô thức} là gì? Nhà triết học khoa học Kuhn [175], khi nghiên cứu cách khoa học tiến hóa thông qua những sự chuyển đổi mô thức, định nghĩa nó là một đóng góp có tính mở, định khung tư duy về một đối tượng nghiên cứu bằng những khái niệm, kết quả và thủ tục tạo cấu trúc cho các thành tựu tương lai. Mặc dù có vẻ xa các mối quan tâm của Khoa học máy tính, định nghĩa này thực sự làm nổi bật sự xuất hiện của một \emph{cấu trúc} (một hệ hình thức) nắm bắt đối tượng bàn luận, và khái niệm về những \emph{thủ tục} (các quy trình) dẫn dắt việc đạt được thành tựu.

\subsection{Các chiều cạnh của mô hình hóa đa mô thức}
Ứng dụng MPM đòi hỏi \emph{mô hình hóa mọi thứ một cách tường minh}, sử dụng \emph{một hoặc nhiều hệ hình thức phù hợp nhất}, ở \emph{một hoặc nhiều mức trừu tượng phù hợp nhất} [266]. Điều này gợi ý rằng có thể hiểu một \emph{mô thức} là một cách sắp xếp các đặc tính trong từng chiều cạnh đã mô tả: các \emph{hệ hình thức} và \emph{mức trừu tượng} trong hoạt động mô hình hóa.

Các hệ hình thức thường có tính đa dụng và khó sử dụng đối với người xây dựng mô hình (người sử dụng trong miền hoặc chuyên gia miền), những người phải bắt đầu bằng việc chọn hệ hình thức thích hợp nhất dựa trên ngữ nghĩa đã biết rõ của nó; chẳng hạn mạng Petri cho luồng công việc và tính đồng thời, hoặc Statecharts để mô tả hệ thống dựa trên sự kiện, cùng những ví dụ khác. Tuy nhiên, việc phải thành thạo ký pháp toán học tạo ra một quá trình học khó khăn. Để giảm nhẹ vấn đề này, các ngôn ngữ chuyên biệt, gọi là \emph{ngôn ngữ mô hình hóa chuyên biệt miền} (Domain-Specific Modelling Languages, DSML), được tạo ra để đơn giản hóa việc biểu đạt ý định đặc tả của người xây dựng mô hình. Các cấu trúc trong những ngôn ngữ này được thiết kế để gần hơn với cách chuyên gia miền quen khái niệm hóa bài toán. Cách tiếp cận có hệ thống để xây dựng ngôn ngữ mô hình hóa mới được gọi là \emph{kỹ nghệ ngôn ngữ mô hình hóa} (Modelling Language Engineering, MLE), và bản thân nó cũng phải tuân theo một quy trình kỹ nghệ [?].

Để xử lý độ phức tạp trong quá trình phát triển hệ thống, ba cách tiếp cận trừu tượng hóa cơ bản thường được kết hợp: \emph{trừu tượng hóa/tinh chỉnh}, \emph{phân rã kiến trúc} và \emph{phân rã khung nhìn}.
\begin{itemize}
\item \textbf{Trừu tượng hóa mô hình (và phép đối ngẫu của nó là tinh chỉnh)} được sử dụng khi tập trung vào một tập \emph{tính chất} cụ thể cần quan tâm. Trong khi trừu tượng hóa có nghĩa là loại bỏ chi tiết không cần thiết, ngược với tinh chỉnh, cùng một tập tính chất đã chọn phải đúng ở cả mô hình trừu tượng lẫn mô hình chi tiết. Kiểm chứng tính chất trên mô hình trừu tượng dự kiến sẽ rẻ hơn (hoặc đơn giản hơn) so với trên mô hình chi tiết. Tuy nhiên, mô hình chi tiết hơn vẫn có những lợi thế vì nó cho phép đánh giá đúng một tập tính chất rộng hơn mà nếu không có nó thì không thể bao quát.
\item \textbf{Phân rã kiến trúc (và phép đối ngẫu của nó là hợp thành thành phần)} được dùng khi có thể chia bài toán thành các phần, mỗi phần có một \emph{giao diện} phù hợp. Cách đóng gói này quy một bài toán về (i) một số bài toán con, mỗi bài toán đòi hỏi thỏa mãn các tính chất riêng và dẫn đến việc thiết kế một thành phần; và (ii) thiết kế một kiến trúc phù hợp để kết nối các thành phần sao cho phép hợp thành thỏa mãn những tính chất yêu cầu ban đầu. Việc chia nhỏ như vậy thường xuất hiện tự nhiên ở một số mức trừu tượng, sử dụng các hệ hình thức phù hợp (có hỗ trợ phân cấp), chẳng hạn nhờ tính cục bộ hoặc liên tục trong miền bài toán/lời giải. Lưu ý rằng mô tả trên là một luồng công việc từ trên xuống, trong đó phân rã các yêu cầu dẫn đến thiết kế thành phần, rồi hợp thành kiến trúc từ các thành phần này. Cũng có thể có một luồng công việc từ dưới lên, kết hợp các thành phần hiện có để thỏa mãn yêu cầu của toàn hệ thống.
% Source: printed 10 / PDF 26
\item \textbf{Phân rã khung nhìn (và phép đối ngẫu của nó là hợp nhất khung nhìn)} được dùng để hỗ trợ cộng tác giữa nhiều bên liên quan, mỗi bên có những mối quan tâm khác nhau. Mỗi góc nhìn cho phép đánh giá một tập tính chất dành riêng cho một bên liên quan. Khi hợp nhất các khung nhìn cụ thể, phép hội của mọi tính chất của các khung nhìn phải đúng. Trong lĩnh vực phần mềm, Tiêu chuẩn IEEE 1471 định nghĩa các quan hệ giữa góc nhìn và hiện thực của chúng là các khung nhìn. Lưu ý rằng các khung nhìn có thể được mô tả bằng những hệ hình thức khác nhau.
\end{itemize}

Một cách kết hợp cụ thể các cách tiếp cận trên dẫn đến Thiết kế hệ thống dựa trên hợp đồng [84]. Thực vậy, các hoạt động mô hình hóa được kết hợp thành \emph{quy trình} (hoặc luồng công việc) liên hệ các hoạt động MPM khác nhau. Quy trình có thể mang tính \emph{mô tả} (descriptive), thể hiện chuỗi hoạt động đã thực hiện cùng các tạo tác liên quan; mang tính \emph{ràng buộc} (proscriptive), bằng cách đặc tả khai báo các ràng buộc đối với những hoạt động được phép và các cách kết hợp chúng; hoặc mang tính \emph{chỉ định} (prescriptive), cho phép thực thi. Quy trình thường được hỗ trợ bằng chuỗi công cụ, trong đó các công cụ khác nhau hỗ trợ những hoạt động khác nhau. Có thể nói một khung MPM nhằm hỗ trợ những người xây dựng công cụ (hoặc siêu công cụ), những người giúp người thực hành suy luận về CPS và xác định hệ hình thức, phép trừu tượng hóa, luồng công việc cùng các phương pháp, kỹ thuật và công cụ hỗ trợ \emph{phù hợp nhất} để thực hiện nhiệm vụ.

Cuối cùng, bằng cách lựa chọn, tổ chức và quản lý ba chiều cạnh trên (hệ hình thức, phép trừu tượng hóa và quy trình), MPM hỗ trợ giao tiếp giữa các chuyên gia, giúp họ nắm bắt rõ hơn bản chất cách xây dựng CPS của mình. Hơn nữa, MPM còn hỗ trợ so sánh chặt chẽ các cách tiếp cận MPM khác nhau dựa trên những thành phần MPM cốt lõi của chúng. Những hệ quả và thách thức mà MPM đặt ra đối với các hệ hình thức và cơ chế trừu tượng hóa cần được thảo luận thêm.

\section{Nền tảng cho MPM4CPS}
Cuốn sách này giới thiệu một tập hợp đại diện các hệ hình thức mô hình hóa, mỗi hệ có một tập đặc tính đặc trưng. Tập hợp này không hề đầy đủ mà nhằm thể hiện sự đa dạng của các đặc tính có trong những hệ hình thức đã được xác lập, thường được hỗ trợ bởi các công cụ có khả năng mở rộng quy mô. Có thể dùng chúng để lựa chọn hệ hình thức phù hợp nhất cho một nhiệm vụ cụ thể; làm điểm khởi đầu để tìm hiểu thêm những hệ hình thức có các đặc tính mong muốn khác (chẳng hạn phân bố không gian như trong ô-tô-mát tế bào hoặc phương trình đạo hàm riêng); hoặc làm cơ sở thiết kế ngôn ngữ mô hình hóa chuyên biệt miền (DSML) để ràng buộc người xây dựng mô hình tối đa vào một miền ứng dụng cụ thể.

Các hệ hình thức giới thiệu trong sách cũng có thể được kết hợp thành những ngôn ngữ ``lai'' khi cần một tổ hợp đặc tính cụ thể mà không có trong một hệ hình thức đơn lẻ. Tuy nhiên, lưu ý rằng một số hệ hình thức trong sách đã mang tính lai theo nghĩa trên. Chẳng hạn, đồ thị liên kết năng lượng thống nhất việc mô hình hóa hệ thống thuộc nhiều miền vật lý bằng cách tập trung vào dòng công suất, còn Modelica kết hợp các đặc tính của hướng đối tượng với mô hình hóa phi nhân quả về mặt tính toán (dựa trên phương trình). Ngôn ngữ phân tích và thiết kế kiến trúc (Architecture Analysis and Design Language, AADL), tập trung vào hệ thống nhúng với kiến trúc làm cốt lõi, tập hợp các góc nhìn khác nhau, nên phù hợp để lập tài liệu, phân tích và tổng hợp mã.

Nội dung được trình bày theo cách từ dưới lên. Trước hết là các hệ hình thức để mô hình hóa thành phần vật lý. Sau đó, các cơ chế đóng gói và tái sử dụng mô tả thành phần CPS được trình bày như phương tiện kiểm soát độ phức tạp của những mô tả lớn. Tiếp theo là các hệ hình thức để phân tích hành vi CPS. Cuối cùng, các hệ hình thức này được kết hợp với nhau thông qua mô tả kiến trúc và quy trình.

\subsection{Mô hình hóa các thành phần vật lý}
Khi mô hình hóa một hệ thống vật lý, quyết định đầu tiên cần đưa ra là liệu các tính chất cần quan tâm của hệ thống có phụ thuộc vào chiều không gian hay không. Chẳng hạn, sự nóng lên của một vật kim loại do dòng điện chạy qua được quyết định bởi tương tác giữa miền vật lý điện và nhiệt. Nó phụ thuộc vào hình học của vật cũng như các tính chất vật liệu của vật, như khối lượng riêng, độ dẫn điện, hằng số điện môi tương đối, nhiệt dung, độ dẫn nhiệt và sự phân bố của chúng trên toàn vật. Động lực học của vật sau đó có thể được mô tả bằng biểu thức toán học của các quan hệ giữa những đại lượng vật lý cần quan tâm. Do phụ thuộc vào tọa độ không gian, việc này đòi hỏi sử dụng các mô hình ``tham số phân bố''. Chúng thường được biểu đạt bằng hệ hình thức phương trình đạo hàm riêng (Partial Differential Equation, PDE).

% Source: printed 11 / PDF 27
Khi các tham số của một vật đủ đồng nhất trên toàn bộ hình học của nó, những tính chất cần quan tâm có thể không phụ thuộc vào chiều không gian. Khi đó, các tham số có thể được gộp trên toàn vật, và vật có thể được quy về phần cốt lõi không có kích thước. Chẳng hạn, một vật rắn trong miền cơ học có khối lượng riêng không đổi trên toàn bộ hình học có thể được quy về một ``chất điểm'' đơn giản. Có thể mô tả động lực học của nó bằng các định luật Newton hoặc cách xây dựng Hamilton hay Lagrange. Những mô hình ``tham số tập trung'' như vậy thường được biểu đạt bằng các hệ hình thức phương trình vi phân thường (Ordinary Differential Equation, ODE) hoặc phương trình vi phân--đại số (Differential Algebraic Equation, DAE).

Các thành phần vật lý của một CPS thường trải rộng trên những miền vật lý khác nhau (điện, cơ học, nhiệt động lực học, thủy lực, v.v.). Hệ hình thức đồ thị liên kết năng lượng được mô tả trong Chương 2 thống nhất những miền này ở mức chi tiết ``tham số tập trung''. Nó nhận ra sự tương đồng giữa các quá trình vật lý thuộc những miền vật lý khác nhau, như tích trữ và tiêu tán năng lượng. Một hệ thống được mô hình hóa dưới dạng đồ thị liên kết năng lượng, kết nối các nút đại diện cho phần tử vật lý. Các nút mã hóa quan hệ giữa những đại lượng vật lý như điện áp và dòng điện, chẳng hạn trong một điện trở. Các cạnh của đồ thị liên kết năng lượng---gọi là \emph{liên kết công suất} (Power Bonds)---biểu thị dòng công suất giữa các nút. Những nút đặc biệt---\emph{nút nối} (junctions)---mã hóa các định luật bảo toàn, là sự tổng quát hóa các định luật dòng điện và điện áp Kirchoff trong miền điện. Chương này giới thiệu một quy trình có hệ thống để mô hình hóa các hệ thống vật lý đa miền. Quy trình bắt đầu từ các Mô hình vật lý lý tưởng hóa và chuyển chúng thành mô hình đồ thị liên kết năng lượng. Các mô hình đồ thị liên kết năng lượng này phi nhân quả về mặt tính toán và có thể được chuyển thành một hệ phương trình vi phân--đại số (DAE). Mô hình phi nhân quả về mặt tính toán gồm các phương trình liên hệ các tín hiệu (các biến, là hàm của thời gian liên tục), mà không chỉ định biến nào đã biết (đầu vào), biến nào chưa biết (đầu ra), cũng như không chỉ định cần giải các phương trình này ra sao (tức là cách tính các đại lượng chưa biết từ các đại lượng đã biết). Những DAE như vậy có thể được biểu diễn bằng các ngôn ngữ mô hình hóa dựa trên phương trình (toán học), chẳng hạn Modelica. Modelica được mô tả trong Chương 3. Chương về đồ thị liên kết năng lượng tiếp đó trình bày cách gán tính nhân quả về mặt tính toán, qua đó chuyển sang một sơ đồ khối nhân quả thời gian liên tục (Continuous-Time Causal Block Diagram, CT-CBD). Sơ đồ khối nhân quả được mô tả trong Chương 4. Gán tính nhân quả trên một mô hình đồ thị liên kết năng lượng có thể đem lại hiểu biết dựa trên vật lý về những sai sót trong mô hình. Sự hỗ trợ ``gỡ lỗi mô hình'' này có được nhờ tính chuyên biệt miền (vật lý) của hệ hình thức đồ thị liên kết năng lượng.

\subsection{Kết nối phần ``vật lý'' với phần ``ảo''}
Các hệ thống thực-ảo được hợp thành từ các thành phần vật lý và tính toán nối mạng. Để thiết kế những hệ thống này theo mô-đun và phân cấp, tối đa hóa khả năng tái sử dụng mô hình và cho phép xây dựng thư viện mô hình, ngôn ngữ Modelica, được mô tả trong Chương 3, kết hợp các đặc tính của hướng đối tượng như đóng gói và kế thừa với mô hình hóa phi nhân quả về mặt tính toán (dựa trên phương trình). Mô hình phi nhân quả về mặt tính toán cho phép người xây dựng mô hình biểu đạt các định luật vật lý cơ bản bằng phương trình toán học. Ngữ nghĩa của Modelica được xác định bằng việc mở rộng các cấu trúc hướng đối tượng như kế thừa, tạo thể hiện của các lớp và làm phẳng cấu trúc phân cấp. Kết quả là một hệ phương trình vi phân--đại số (lai). Với mỗi bối cảnh thực nghiệm mô phỏng cụ thể, trình biên dịch Modelica có thể gán tính nhân quả về mặt tính toán. Điều này thực chất sinh ra một mô hình trong hệ hình thức sơ đồ khối nhân quả thời gian liên tục. Hầu hết trình biên dịch Modelica sẽ tiếp tục rời rạc hóa (theo thời gian) các mô hình này, bằng phương pháp ký hiệu thông qua ``tích phân nội tuyến'' (inline integration), hoặc gọi các bộ giải số bên ngoài. Cuối cùng, việc này tạo ra mã mô phỏng thực thi được. Lưu ý rằng có thể tạo một thư viện Modelica với các thành phần đồ thị liên kết năng lượng. Khi đó, quy trình gán tính nhân quả của đồ thị liên kết năng lượng không được sử dụng. Thay vào đó, phép gán tính nhân quả của Modelica được áp dụng cho toàn mô hình, kể cả những phần không thuộc đồ thị liên kết năng lượng. Thông qua đặc tả hàm bằng mã, Modelica còn cho phép biểu diễn mã hướng đối tượng truyền thống. Chính sự kết hợp giữa mã, phương trình và các cấu trúc lai như \texttt{when} (cho phép đưa vào những sự kiện rời rạc, gọi là ``sự kiện trạng thái'', dựa trên các điều kiện đối với hành vi liên tục như vượt qua một giá trị ngưỡng) khiến Modelica phù hợp để xây dựng mô hình CPS, bao quát các phần vật lý, mạng và tính toán của chúng.

Như đã đề cập, các mô hình khai báo, phi nhân quả về mặt tính toán cuối cùng cần được biến đổi sang dạng nhân quả để có thể giải bằng tính toán. Trong Chương 4, một họ các hệ hình thức sơ đồ khối nhân quả (Causal Block Diagram, CBD) được giới thiệu. Sơ đồ khối nhân quả gồm một mạng các khối tính toán. Mỗi khối đặc tả quan hệ nhân quả về mặt tính toán giữa các tín hiệu đầu vào và đầu ra của nó. Mạng sơ đồ khối đặc tả cách các đầu ra của một khối được kết nối với đầu vào của những khối khác. Kết nối này biểu thị rằng giá trị tại các cổng đầu ra và đầu vào được nối với nhau phải luôn bằng nhau. Ba biến thể CBD được xây dựng
% Source: printed 12 / PDF 28
dần dần. Hệ hình thức sơ đồ khối nhân quả đại số (Algebraic Causal Block Diagram, ALG-CBD) không có khái niệm thời gian: giá trị được truyền qua CBD theo một ``lịch'' tính toán (tức thứ tự gọi các phép tính của khối), suy ra từ cấu trúc phụ thuộc được mã hóa trong mạng sơ đồ khối. Cần đặc biệt chú ý phát hiện và giải đúng các chu trình phụ thuộc, được gọi là ``vòng lặp đại số''. Sau đó, một khái niệm thời gian rời rạc (số tự nhiên) cùng một khối trễ/bộ nhớ được bổ sung, tạo ra sơ đồ khối nhân quả thời gian rời rạc (Discrete-Time Causal Block Diagram, DT-CBD). Chúng tương đương với mô hình luồng dữ liệu đồng bộ (Synchronous Data Flow, SDF). Cuối cùng, số thực được đưa vào làm cơ sở thời gian để tạo sơ đồ khối nhân quả thời gian liên tục (CT-CBD). Những mô hình này có cùng sức biểu đạt với các phương trình toán học và cần được rời rạc hóa để có thể giải bằng tính toán. Các kỹ thuật rời rạc hóa số được dùng để chuyển một CT-CBD thành một DT-CBD.

Trong nhiều trường hợp, hợp lý hơn là trừu tượng hóa bỏ nhiều chi tiết hành vi của hệ thống và chỉ tập trung vào các ``sự kiện'' liên quan. Các phép trừu tượng hóa sự kiện rời rạc như vậy xem hệ thống là thay đổi trạng thái bên trong, hoặc để phản ứng với sự kiện đầu vào, hoặc tự chủ thay đổi trạng thái sau một khoảng thời gian nhất định (do một ``sự kiện thời gian'' bên trong), và có thể tạo ra sự kiện đầu ra tại một số thời điểm. Vì chỉ các sự kiện mới làm thay đổi trạng thái hệ thống, và giữa những lần xuất hiện sự kiện ta giả định không có điều gì liên quan xảy ra, sự tiến triển của trạng thái theo thời gian là hằng từng đoạn. Khác với các hệ hình thức thời gian rời rạc (Discrete-Time, DT), trong các hệ hình thức sự kiện rời rạc (Discrete-Event, DE), thời gian tiến bằng các bước nhảy từ sự kiện liên quan này sang sự kiện liên quan tiếp theo. Một lợi thế của phép trừu tượng hóa DE là hiệu năng: trình mô phỏng bước thẳng đến thời điểm tiếp theo có sự kiện, trong khi trình mô phỏng DT phải tiến qua thời gian bằng những bước tăng cố định dù không có gì đáng chú ý xảy ra (tức trạng thái không đổi). Phép trừu tượng hóa DE thường được dùng để nghiên cứu cạnh tranh giữa các quá trình khác nhau đối với tài nguyên dùng chung. Nếu tài nguyên bị giới hạn, điều này tất yếu dẫn đến xếp hàng. Vì vậy, phép trừu tượng hóa này hữu ích cho phân tích hiệu năng dựa trên mô phỏng. Mức sử dụng tài nguyên, thời gian hoàn thành một hoạt động, phân bố độ dài hàng đợi và thời gian xếp hàng đều là ví dụ về các thước đo hiệu năng điển hình thu được từ mô phỏng sự kiện rời rạc. Lưu ý rằng các hệ hình thức DE thường tất định. Bằng cách dùng các phân bố cho giá trị tham số như thời gian giữa các lần đến (Inter Arrival Time), thay cho những giá trị duy nhất, và sử dụng mô phỏng Monte Carlo, ta thu được các phân bố của thước đo hiệu năng. Nhờ đó, ta có các mô phỏng ngẫu nhiên có thể lặp lại (vì các bộ sinh số giả ngẫu nhiên---vốn tất định---được dùng để lấy mẫu từ các phân bố). Qua nhiều năm, nhiều hệ hình thức DE đã được phát triển. Hệ hình thức Đặc tả sự kiện rời rạc (Discrete EVent Specification, DEVS), được mô tả trong Chương 5, là một hệ hình thức DE đủ cơ bản và đủ sức biểu đạt để đóng vai trò ``hợp ngữ'' DE: mô hình trong mọi hệ hình thức mô phỏng DE có thể được ánh xạ sang một mô hình DEVS tương đương. Vì vậy, có thể dùng nó để kết nối theo kiến trúc các mô hình thuộc những hệ hình thức khác nhau, bằng cách trước hết ánh xạ mọi thành phần sang DEVS. Kiến trúc thu được chỉ chứa các thành phần DEVS. Vì DEVS có tính mô-đun và hỗ trợ hợp thành kiến trúc phân cấp, mô hình thu được có ý nghĩa chính xác. Khả năng hỗ trợ phân cấp của DEVS khiến nó phù hợp để xây dựng thư viện mô hình, rồi xây dựng những mô hình rất phức tạp.

Một cách khác để kết hợp các ô-tô-mát trạng thái theo mô-đun nằm trong hệ hình thức Statecharts, được mô tả trong Chương 6. Hệ hình thức Statecharts gồm các cấu trúc phân cấp ô-tô-mát trạng thái, phép hợp thành song song các ô-tô-mát này, một khái niệm thời gian và cơ chế quảng bá sự kiện. Sự phổ biến của Statecharts một phần nhờ ký pháp trực quan dễ hiểu của nó. Mục đích chính của Statecharts không chỉ là mô phỏng mô hình, mà còn tổng hợp từ chúng phần mềm và/hoặc phần cứng tự chủ, có yếu tố thời gian và phản ứng.

Khi hợp thành song song các ô-tô-mát trạng thái, có thể có nhiều cách xen kẽ. Một lựa chọn là chọn một cách xen kẽ duy nhất, dẫn đến một vết hành vi duy nhất và tất định. Đây là cách các hệ hình thức DEVS và Statecharts thực hiện. Để mô hình hóa tính đồng thời thực sự, việc tuần tự hóa nhân tạo này không phải lúc nào cũng phù hợp. Thay vào đó, nên cho phép lựa chọn không tất định. Điều này dẫn đến một tập hợp các vết hành vi có thể có, thay vì một vết hành vi duy nhất. Có thể tóm lược tập hợp vết này bằng một biểu diễn gọn dưới dạng đồ thị khả đạt trạng thái. Sau đó có thể kiểm tra việc thỏa mãn những tính chất cần quan tâm trên tập hợp vết. Một ví dụ là khả năng đạt đến một trạng thái không mong muốn nào đó. Những tính chất này cũng được biểu đạt bằng một ngôn ngữ tính chất phù hợp. Các ngôn ngữ không tất định như mô tả trên thường có một khái niệm thời gian yếu: cơ sở thời gian không dùng số tự nhiên hay số thực, mà chỉ dùng một thứ tự (bộ phận) của các thời điểm sự kiện. Trọng tâm là tính đồng thời và đồng bộ hóa. Các hệ hình thức này chủ yếu được dùng để phân tích tính chất trên mọi hành vi có thể có của hệ thống, thay vì mô phỏng hoặc tổng hợp. Điều này khiến chúng phù hợp, chẳng hạn, cho phân tích an toàn. Một hệ hình thức như vậy là mạng Petri, được mô tả trong Chương 7. Mạng Petri mã hóa trạng thái dưới dạng một ``đánh dấu'' (marking), một vectơ $n$ chiều gồm các số tự nhiên. Mỗi phần tử của vectơ tương ứng với số lượng ``thẻ'' (token) trong một vị trí (Place) của mạng Petri. Sự tiến triển của trạng thái được mã hóa trong một đồ thị mạng Petri, ngoài các vị trí còn chứa những chuyển tiếp và cung. Nhờ sử dụng số tự nhiên, số trạng thái có thể có có thể vô hạn (nhưng đếm được). Mạng Petri là một hệ hình thức đơn giản, giống như DEVS, thường được dùng làm miền ngữ nghĩa chung để ánh xạ nhiều hệ hình thức khác,
% Source: printed 13 / PDF 29
thường là chuyên biệt miền, vào đó. Chương này cũng minh họa cách kết hợp mạng Petri với những hệ hình thức khác. Đặc biệt, chương giới thiệu đồng mô phỏng mạng Petri với các Functional Mockup Unit (mã hóa mô hình liên tục đã được rời rạc hóa). Điều này thực chất dẫn đến các mô hình lai không tất định.

\subsection{Công cụ hỗ trợ cho MPM4CPS}
Những hệ thống kỹ nghệ phức tạp gồm các thành phần không đồng nhất được bố trí trong một kiến trúc. Hơn nữa, nhiều góc nhìn về cùng hệ thống có thể được quan tâm, và cuối cùng (một phần của) mô hình hệ thống cần được ``triển khai'' trên một kiến trúc phần mềm/phần cứng, thường là kiến trúc nhúng. Ngôn ngữ phân tích và thiết kế kiến trúc (AADL), được mô tả trong Chương 8.1, trước hết là một ngôn ngữ mô tả kiến trúc. Nó cho phép mô tả tổng thể hệ thống và môi trường mà hệ thống sẽ hoạt động trong đó. Từ mô tả này, có thể sinh những mô hình khác thuộc các hệ hình thức khác, như những hệ hình thức được mô tả trong sách. Có thể tiếp tục bổ sung các mô hình đó để nghiên cứu các khía cạnh khác nhau của hệ thống, điều thiết yếu cho tối ưu hóa, kiểm chứng và thẩm định. Sau phần giới thiệu ngắn về các ngôn ngữ mô tả kiến trúc (Architecture Description Languages, ADL) và vai trò của chúng trong MPM4CPS, AADL được trình bày và cách sử dụng được minh họa qua mô hình hóa, phân tích và sinh mã cho một robot Lego Mindstorm đơn giản để vận chuyển vật trong nhà kho. Quy trình thiết kế đơn giản từ trên xuống, lấy kiến trúc làm trung tâm, được áp dụng: bắt đầu bằng việc nắm bắt mục tiêu của các bên liên quan và yêu cầu hệ thống; tiếp đến thiết kế hệ thống, phân tích và kiểm chứng thiết kế; cuối cùng là tổng hợp mã tự động.

Từ những điều trên có thể thấy rằng mô hình hóa hệ thống phức tạp không chỉ bao gồm những phép trừu tượng hóa, kiến trúc và khung nhìn khác nhau, được mô hình hóa bằng các hệ hình thức khác nhau, mà còn bao gồm các luồng công việc phát triển phức tạp. Chương 9 xem xét chủ đề mô hình hóa quy trình (luồng công việc). Các quy trình phát triển đồng thời, thường phức tạp, được mô hình hóa dưới dạng Mô hình quy trình (Process Model, PM), được xây dựng từ những hoạt động nguyên thủy nhận đầu vào là các tạo tác mô hình hóa và tạo ra đầu ra là các tạo tác mô hình hóa đã sửa đổi hoặc mới. Các hoạt động có thể cần nguồn lực con người hoặc máy tính, có thể gây ra trì hoãn như mô tả trong chương về DEVS. Vì các tạo tác mà hoạt động thao tác là mô hình thuộc nhiều hệ hình thức khác nhau, việc ``gán kiểu'' cho chúng bằng những hệ hình thức phù hợp là hợp lý. Để biểu diễn các hệ hình thức được sử dụng và cho thấy quan hệ giữa chúng, Đồ thị biến đổi hệ hình thức (Formalism Transformation Graph, FTG) được giới thiệu. FTG+PM, kết hợp FTG với PM, cho phép đặc trưng hóa bản chất của các mẫu giải pháp mô hình hóa đa mô thức cho những bài toán phát triển CPS. Một lợi thế của biểu diễn FTG+PM tường minh là có thể dùng nó làm cơ sở tổng hợp các công cụ MPM.

\section{Tổng kết}
Lĩnh vực CPS chịu ảnh hưởng bởi độ phức tạp của những cách tiếp cận, quy trình và ngôn ngữ mô hình hóa khác nhau. Thực vậy, lĩnh vực này được biết đến với bản chất đa ngành; vì không có một cách tiếp cận mô hình hóa duy nhất được chấp nhận rộng rãi, mô hình hóa đa mô thức được đề xuất làm giải pháp. Tuy nhiên, cho đến nay, tài liệu vẫn thiếu một nguồn tham khảo vững chắc giới thiệu những kỹ thuật mô hình hóa và phân tích CPS khác nhau hướng đến một cách tiếp cận đa mô thức.

Chương đầu tiên trình bày động cơ của mô hình hóa đa mô thức cho CPS và nhu cầu về nền tảng rõ ràng của nó. Bắt đầu bằng giới thiệu khái niệm Hệ thống, chương trình bày các lớp hệ thống khác nhau, bàn luận đặc điểm của chúng, rồi suy ra một định nghĩa CPS. Vì CPS khó xây dựng, chúng tôi dành một phần chương để thảo luận và đi sâu vào những nguồn gây ra độ phức tạp thường gặp. Nhìn chung, thiết kế một hệ thống có nghĩa là phải sử dụng nhiều loại phép trừu tượng hóa để mô tả các tính chất và mối quan tâm khác nhau của hệ thống, bằng các ngôn ngữ (với các hệ hình thức và mô thức mô hình hóa hiện có) và quy trình để tránh độ phức tạp không cần thiết. Sau đó, chương định nghĩa mô hình là gì và việc sử dụng nhiều hệ hình thức khi thực hiện mô hình hóa đa mô thức có nghĩa gì. Cuối cùng, cấu trúc của cuốn sách được mô tả, giải thích thêm cách các hệ hình thức và kỹ thuật được trình bày gắn kết với nhau.

% Source: printed 14 / PDF 30
\paragraph*{Truy cập mở}
Chương này được cấp phép theo các điều khoản của Giấy phép quốc tế Creative Commons Ghi công 4.0 (\url{http://creativecommons.org/licenses/by/4.0/}), cho phép sử dụng, chia sẻ, chuyển thể, phân phối và sao chép trên mọi phương tiện hoặc định dạng, miễn là ghi công phù hợp cho tác giả gốc và nguồn, cung cấp liên kết đến giấy phép Creative Commons và cho biết nếu đã thực hiện thay đổi.

Các hình ảnh hoặc tư liệu bên thứ ba khác trong chương này được bao gồm trong giấy phép Creative Commons của chương, trừ khi có chỉ dẫn khác trong dòng ghi công của tư liệu đó. Nếu tư liệu không nằm trong giấy phép Creative Commons của chương và mục đích sử dụng của bạn không được pháp luật cho phép hoặc vượt quá phạm vi sử dụng được cho phép, bạn cần xin phép trực tiếp từ chủ sở hữu bản quyền.
